\section{Testing, done properly}

The prompt says plainly: \emph{``The quality of your test suite is part of the
evaluation \dots\ we are looking for evidence of thoughtful test design.''} That
means tests are graded like production code. This section gives you a
\emph{theory} of how to test this problem, not just a list --- so you can justify
every test and add the ones most candidates forget: an \textbf{oracle}, a
\textbf{property-based} layer, \textbf{consistency} tests, and a
\textbf{performance} test.

\subsection{The four layers of a convincing suite}

\begin{center}
\begin{tikzpicture}[font=\footnotesize,>=Latex]
  \tikzstyle{lay}=[draw,rounded corners,minimum width=12.5cm,minimum height=0.95cm,align=left,inner sep=6pt]
  \node[lay,fill=Buy!10] (l1) {\textbf{1. Example tests} --- the three README cases; the ground truth everyone must pass.};
  \node[lay,fill=Accent!12,below=3mm of l1] (l2) {\textbf{2. Edge / boundary tests} --- empty, one company, only buys, \texttt{minQty} boundary, overflow.};
  \node[lay,fill=Gold!14,below=3mm of l2] (l3) {\textbf{3. Consistency / metamorphic} --- add\,$\to$\,cancel\,$\to$\,re-query; order independence; index integrity.};
  \node[lay,fill=Sell!12,below=3mm of l3] (l4) {\textbf{4. Property-based + oracle} --- random inputs checked against an independent reference (max-flow).};
\end{tikzpicture}
\end{center}

\begin{intu}
Layers 1--2 prove ``it works on the cases I \emph{thought} of''. Layers 3--4 prove
``it works on cases I did \emph{not} think of'' --- which is the whole point of
testing. Interviewers notice immediately when you climb past layer 2, because
almost nobody does.
\end{intu}

\subsection{Layer 1 --- the example tests (necessary, not sufficient)}

The three README examples ($2700$, and the Example~2 / Example~3 tables) are
non-negotiable. Your starter suite already has them, e.g.:

\begin{lstlisting}
EXPECT_EQ(cache.getMatchingSizeForSecurity("SecId2"), 2700);
EXPECT_EQ(cache.getMatchingSizeForSecurity("SecId1"), 0);   // same-company
EXPECT_EQ(cache.getMatchingSizeForSecurity("SecId3"), 0);   // only buys
\end{lstlisting}

\begin{pitfall}
Passing the three examples proves \emph{almost nothing} on its own --- a wrong
\code{min(B,S)} implementation passes two of the three SecIds in Example~1! You
must add cases that \emph{distinguish} the right formula from the plausible wrong
ones. That is what the next layers are for.
\end{pitfall}

\subsection{Layer 2 --- edges chosen to kill specific bugs}

Every edge case should be traceable to a bug it would catch. Design them from the
formula $M=\min(B,S,\,B+S-\max_c(b_c+s_c))$: each term suggests a test.

\begin{center}
\renewcommand{\arraystretch}{1.3}\footnotesize
\begin{tabular}{@{}p{4.3cm}p{5.5cm}p{4.0cm}@{}}
\toprule
\textbf{Test} & \textbf{Input} & \textbf{Catches} \\
\midrule
one company only & \Buy{}500(A), \Sell{}500(A) $\to$ \textbf{0} & forgetting the company rule \\
one cross match & \Buy{}100(A), \Sell{}100(A), \Sell{}100(B) $\to$ \textbf{100} & counting own-company sells \\
\textbf{dominant company} & \Buy{}1000(A), \Sell{}900(A), \Buy{}100(B), \Sell{}50(C) $\to$ \textbf{150} & using naive $\min(B,S){=}950$ \\
only buys / only sells & all one side $\to$ \textbf{0} & missing empty-side guard \\
partial match & \Buy{}100, \Sell{}300 (diff co) $\to$ \textbf{100} & wrong side cap \\
\code{minQty} boundary & cancel with qty \emph{exactly} $=$ minQty & \code{>} vs \code{>=} off-by-one \\
unknown security / id / user & any lookup $\to$ no-op, no throw & missing graceful guard \\
overflow & many orders of qty near \code{UINT\_MAX} & 32-bit accumulator \\
\bottomrule
\end{tabular}
\end{center}

\begin{keybox}[The star edge case: dominant company]
\Buy{}1000(A), \Sell{}900(A), \Buy{}100(B), \Sell{}50(C). Here $B{=}1100$,
$S{=}950$, so a naive implementation returns $\min(1100,950){=}\mathbf{950}$. The
\emph{correct} answer is $\mathbf{150}$ (company A's huge position can only trade
with B and C's tiny opposite side). This single test separates a correct solution
from the most common wrong one. Put it in your suite and \emph{mention it}.
\end{keybox}

\subsection{Layer 3 --- consistency and metamorphic tests}

These check \emph{relationships} between operations rather than fixed answers.
They are cheap to write and catch the desync bugs from \S5.

\begin{itemize}[leftmargin=1.5em]
  \item \textbf{Add--cancel round trip.} After \code{addOrder} then
        \code{cancelOrder} of the same id, the cache must be byte-for-byte back to
        its earlier state: same \code{getAllOrders} size, \emph{and} same matching
        numbers. This catches aggregates that are updated on add but not on cancel.
  \item \textbf{Order independence (commutativity).} Insert the same orders in a
        different order; every \code{getMatchingSizeForSecurity} must be identical.
        Verified true in theory (matching depends only on the multiset of orders).
        A failure means hidden order-dependent state. \emph{(0 violations across
        50{,}000 random shuffles in my check.)}
  \item \textbf{Self-pair monotonicity (metamorphic).} Adding a brand-new
        company's matching buy+sell pair can never \emph{decrease} any security's
        match size. \emph{(0 violations in 50{,}000 trials.)}
  \item \textbf{Index integrity after bulk cancels.} After
        \code{cancelOrdersForUser}, no leftover order in \code{getAllOrders}
        belongs to that user, and querying that user again is a clean no-op.
\end{itemize}

\begin{lstlisting}
TEST_F(OrderCacheTest, AddCancelRoundTripRestoresMatching) {
    seedExample1(cache);                       // some fixed scenario
    auto before = cache.getMatchingSizeForSecurity("SecId2");
    cache.addOrder({"TMP","SecId2","Sell",123,"UX","CompZ"});
    cache.cancelOrder("TMP");
    EXPECT_EQ(cache.getMatchingSizeForSecurity("SecId2"), before); // aggregate healed
}
\end{lstlisting}

\subsection{Layer 4 --- property-based testing with an oracle}

This is the layer that turns ``I hope it's right'' into ``I \emph{proved} it's
right on 100k cases''. The idea, from QuickCheck-style testing:

\begin{center}
\begin{tikzpicture}[font=\footnotesize,>=Latex,node distance=8mm,every node/.style={align=center}]
  \node[draw,rounded corners,fill=Soft] (gen) {random\\orders};
  \node[draw,rounded corners,fill=Buy!10,right=1.4cm of gen] (impl) {your\\OrderCache};
  \node[draw,rounded corners,fill=Accent!12,right=1.4cm of impl] (orc) {independent\\\textbf{oracle}};
  \node[draw,rounded corners,fill=Gold!14,right=1.4cm of orc] (cmp) {ASSERT\\equal?};
  \draw[->](gen)--(impl);\draw[->](gen)-|(orc);\draw[->](impl)--(cmp);\draw[->](orc)|-(cmp);
  \node[below=1.2cm of impl,align=center,font=\itshape\color{gray}]
    {two \emph{different} computations of the same answer must agree};
\end{tikzpicture}
\end{center}

The oracle must be written a \emph{different way} than the implementation so they
don't share a bug. Two good oracles here:

\begin{enumerate}[leftmargin=1.5em]
  \item \textbf{Brute max-flow} (Edmonds--Karp on the tiny company graph). This is
        the gold standard --- it computes the true optimum from first principles.
        It is what I used to validate both closed forms (0 mismatches on
        $200{,}000$ instances). Slow, but you only run it on small random cases.
  \item \textbf{The \emph{other} closed form.} Implement Form~A in the cache and
        Form~B in the test (or vice-versa); they must agree. Cheap, and it caught
        nothing only because both were already correct --- but it's a great
        regression guard.
\end{enumerate}

\begin{lstlisting}
TEST_F(OrderCacheTest, RandomizedAgainstMaxFlowOracle) {
    std::mt19937 rng(12345);                       // fixed seed = reproducible
    for (int iter = 0; iter < 2000; ++iter) {
        OrderCache c;
        auto scenario = randomOrders(rng);         // small: few companies, qtys
        for (auto& o : scenario) c.addOrder(o);
        for (auto& sec : securitiesIn(scenario))
            EXPECT_EQ(c.getMatchingSizeForSecurity(sec),
                      maxFlowOracle(scenario, sec)) // independent reference
                << "seed-derived scenario iter=" << iter;
    }
}
\end{lstlisting}

\begin{intu}
Why a \emph{fixed seed}? Reproducibility. A property test that fails must fail the
\emph{same way} every run so you can debug it. Print enough state
(\code{<< scenario}) that a red test hands you the exact counter-example.
\end{intu}

\subsection{Performance testing (the ``1 million orders'' clause)}

The README demands efficiency at scale. Add \emph{one} test that proves it:

\begin{lstlisting}
TEST_F(OrderCacheTest, MillionOrdersStaysFast) {
    OrderCache c;
    for (int i = 0; i < 1'000'000; ++i)
        c.addOrder(makeOrder(i));                  // O(1) each
    auto t0 = std::chrono::steady_clock::now();
    volatile auto m = c.getMatchingSizeForSecurity("HotSec");
    auto ms = std::chrono::duration_cast<std::chrono::milliseconds>(
                  std::chrono::steady_clock::now() - t0).count();
    EXPECT_LT(ms, 50);      // aggregate query must be ~O(#companies), not O(n)
}
\end{lstlisting}

\begin{pitfall}
A wrong (scan-every-order) matching implementation \emph{still passes correctness}
but blows this timing test. That is the point: the performance test is what proves
your \emph{aggregate} design actually pays off. Keep the bound generous (CI
machines vary) but tight enough to fail an $O(n)$ query.
\end{pitfall}

\subsection{Coverage checklist (map every method to a test)}

\begin{center}
\renewcommand{\arraystretch}{1.25}
\begin{tabular}{@{}ll@{}}
\toprule
\textbf{Method / concern} & \textbf{Covered by} \\
\midrule
addOrder (valid / duplicate / invalid) & AddAndRetrieve, DuplicateIgnored, InvalidIgnored \\
cancelOrder (present / absent) & CancelOrder, CancelNonexistentNoOp \\
cancelOrdersForUser & CancelForUser + index-integrity check \\
cancelOrdersForSecIdWithMinimumQty & boundary test asserting \emph{which ids remain} \\
getMatchingSizeForSecurity & all 3 README + dominant-company + oracle \\
getAllOrders & every test that inspects contents \\
concurrency (if implemented) & multi-thread add/cancel + \code{-fsanitize=thread} \\
performance & MillionOrdersStaysFast \\
\bottomrule
\end{tabular}
\end{center}

\begin{pitfall}
Your starter \code{CancelOrdersForSecIdWithMinimumQty} test asserts only the
\emph{count} (\code{size()==2}) and leaves the identities in a comment. Upgrade it
to assert the \emph{exact surviving ids} (\code{O1} and \code{O4}) --- otherwise a
bug that deletes the wrong two orders still passes.
\end{pitfall}

\subsection{Tooling that makes tests worth more}

\begin{itemize}[leftmargin=1.5em,label=\textcolor{Accent}{\faTools}]
  \item \textbf{Sanitizers.} Build the suite with
        \code{-fsanitize=address,undefined} --- iterator invalidation and integer
        overflow become \emph{hard failures}, not silent wrong answers. Add
        \code{-fsanitize=thread} if you implement locking.
  \item \textbf{Fixed seeds + reproducible RNG} for every randomized test.
  \item \textbf{One assertion idea per test} where possible, with a descriptive
        name (\code{DominantCompanyNotNaiveMin}) so a failure reads like a sentence.
  \item \textbf{A tiny \code{seedExampleN} helper} so scenarios are written once
        and reused, keeping tests readable.
\end{itemize}

\begin{say}
``My suite has four layers: the README examples, edge cases each targeting a
specific bug --- the key one being a \emph{dominant-company} case where naive
$\min(B,S)$ over-counts --- consistency tests like add/cancel round-trips and
order-independence, and a property-based layer that checks thousands of random
scenarios against an independent max-flow oracle with a fixed seed. I also add a
million-order timing test to prove the matching query is $O(\#\text{companies})$,
and I run everything under ASan/UBSan so overflow and iterator bugs fail loudly.''
\end{say}
